\chapter{A Primeira Conversa: do Pedido Vago ao Resultado Útil}
\label{ch:primeira-conversa}

\ficcao O gerente abre um chat e escreve: «Resolve o problema das meias desaparecidas.» Pode receber uma lista de hipóteses, uma campanha de marketing ou uma teoria sobre duendes. Nenhuma dessas respostas demonstra onde está uma meia. O problema começa no pedido: não define uma tarefa que a ferramenta possa executar e verificar.

Chamamos \emph{prompt} ao pedido e às instruções que escrevemos para orientar a resposta. Aqui chega um objetivo pequeno. Queremos preparar um rascunho para um cliente, com a informação que já temos.

\section{Dar ao pedido uma finalidade}
\ficcao A mensagem do cliente é: «Entreguei seis pares na segunda-feira e recebi cinco. Podem confirmar o que aconteceu?» Não temos acesso ao registo de recolha. Sabemos apenas que a política fictícia pede o número do talão antes de investigar.

Podemos copiar este pedido para um chat.

\begin{lstlisting}
Prepara um rascunho de resposta em portugues de Portugal.
O cliente diz que entregou seis pares e recebeu cinco.
Nao temos o registo de recolha nem sabemos a causa.
A politica pede o numero do talao para investigar.

A resposta deve reconhecer o pedido, pedir esse numero
e explicar o proximo passo. Usa ate quatro frases.
Nao prometas reembolso, prazo nem resultado da investigacao.
\end{lstlisting}

O primeiro parágrafo diz o que sabemos. O segundo diz o que fazer com isso. Qualquer redação que reconheça a reclamação e peça o talão sem atribuir uma causa pode servir de rascunho.

\begin{keyideabox}{Especificar é reduzir adivinhas}
Pergunta: o que quero obter, com que informação, sob que limites e como vou verificar? O papel do modelo é opcional; esses quatro elementos são mais importantes do que começar por «és o maior especialista do mundo».
\end{keyideabox}

\section{Uma resposta ilustrativa e uma verificação real}
Uma resposta ilustrativa, escrita para o exercício, é «Obrigado pelo aviso. Para verificarmos a diferença entre a entrega e a recolha, pode indicar o número do talão? Com esse número, a equipa poderá consultar o registo e dar seguimento à reclamação.» Usa-a como referência para comparar a resposta que obtiveres.

Assinala quatro verificações. A resposta reconhece o pedido? Pede o talão? Mantém a causa em aberto? Evita prometer um prazo ou uma compensação? «Lamentamos o duende ter fugido com a sua meia» falha o terceiro critério, mesmo que tenha melhor sentido de humor.

Também podemos distinguir factos atribuídos de factos confirmados. «O cliente refere que faltou um par» descreve a mensagem. «A lavandaria perdeu um par» afirma uma conclusão ainda não investigada. Uma palavra muda a responsabilidade da frase.

\section{Melhorar uma coisa de cada vez}
Se a resposta for demasiado longa, especifica o limite e confirma-o. Se inventar uma causa, torna explícita a informação em falta e testa outra mensagem. Se o formato variar, mostra um exemplo. Evita acrescentar quinze regras sem saber que falha cada uma tenta resolver.

Para comparar duas versões, guarda a mesma entrada e usa sessões independentes, quando a interface o permite. Caso contrário, uma resposta anterior pode influenciar a seguinte. Não confundas mudança de modelo com efeito do prompt. Regista a ferramenta, data, versão identificável e configuração disponível.

Num chat não controlas necessariamente todos os detalhes: pode haver instruções da plataforma, pesquisa automática, memória ou seleção de ferramentas. Anota o que consegues observar. «Não sei qual era a configuração interna» é uma descrição honesta da experiência.

\begin{promptlabbox}{A primeira reclamação da Lavandaria}
\textbf{Objetivo e contexto:} comparar um pedido vago com uma tarefa delimitada. Usa a mensagem dos seis pares e a política fictícia deste capítulo. Não forneças dados de clientes reais.

\textbf{Instruções:} em duas sessões, fornece a mesma mensagem e a mesma política. Numa, acrescenta «responde ao cliente». Na outra, acrescenta as instruções do exemplo. Guarda as duas saídas. Repete com «não tenho o talão» e com «quero uma indemnização já».

\textbf{Resultado esperado:} o segundo pedido torna observáveis os critérios da resposta, sem garantir acerto. Para a falta de talão, deve identificar o bloqueio ou pedir um procedimento autorizado; não inventar uma exceção.

\textbf{Validação:} usa os quatro critérios acima, marcando cumpre/não cumpre e copiando o trecho que justifica a avaliação. Se quiseres quantificar variação, repete cada condição e guarda todas as respostas, incluindo falhas.

\textbf{Falhas possíveis:} avaliar só a simpatia, escolher apenas a melhor tentativa, comparar entradas diferentes ou tratar um resultado ilustrativo como resposta medida.
\end{promptlabbox}

\section{Do rascunho à ação}
Enquanto estás a experimentar, podes rever uma frase antes de a usar. Numa aplicação automática, o texto pode ser enviado a milhares de pessoas. A diferença de escala altera o que precisamos de testar.

Antes de responder a outra reclamação, guarda a mensagem, o pedido e o resultado das quatro verificações. Quando trocares uma frase, esse registo permite ver se corrigiste a falha. É o ponto de partida do próximo capítulo.
